\documentclass[11pt,a4paper]{article}

% A self-contained mathematical working note.  No external .bib or .sty files.
\usepackage[T1]{fontenc}
\usepackage{lmodern}
\usepackage{microtype}
\usepackage{geometry}
\geometry{margin=30mm,marginparwidth=18mm}
\usepackage{amsmath,amssymb,amsthm,mathtools}
\usepackage{aliascnt}
\usepackage[dvipsnames]{xcolor}
\usepackage[colorlinks=true,linkcolor=MidnightBlue,citecolor=MidnightBlue,
            urlcolor=MidnightBlue,bookmarksnumbered=true]{hyperref}
\usepackage[nameinlink,capitalise,noabbrev]{cleveref}

% Shared numbering with distinct hyperlinked reference names.
\newtheorem{theorem}{Theorem}[section]
\newaliascnt{proposition}{theorem}
\newtheorem{proposition}[proposition]{Proposition}
\aliascntresetthe{proposition}
\newaliascnt{lemma}{theorem}
\newtheorem{lemma}[lemma]{Lemma}
\aliascntresetthe{lemma}
\newaliascnt{corollary}{theorem}
\newtheorem{corollary}[corollary]{Corollary}
\aliascntresetthe{corollary}
\theoremstyle{definition}
\newaliascnt{definition}{theorem}
\newtheorem{definition}[definition]{Definition}
\aliascntresetthe{definition}
\newaliascnt{example}{theorem}
\newtheorem{example}[example]{Example}
\aliascntresetthe{example}
\theoremstyle{remark}
\newaliascnt{remark}{theorem}
\newtheorem{remark}[remark]{Remark}
\aliascntresetthe{remark}

\newcommand{\cL}{\mathcal{L}}
\newcommand{\cK}{\mathcal{K}}
\newcommand{\weakto}{\rightharpoonup}
\newcommand{\N}{\mathbb{N}}
\newcommand{\ellone}{\ell^{1}}

\title{Compact Operators and Weak Convergence\\
       \large A Working Note on Sequential Regularity}
\author{}
\date{October 2026 \quad (illustrative note)}

\begin{document}
\maketitle

\begin{abstract}
The purpose of this note is to keep together a small chain of observations
about compact operators, weak convergence, and the role of reflexivity.
The exposition moves between precise statements, proofs, a counterexample,
and questions suggested by the comparison.  The references and cross-links
are part of the argument rather than a separate catalogue of facts.
\end{abstract}

\tableofcontents

\section{The distinction to keep in mind}\label{sec:setting}

Let $X$ and $Y$ be Banach spaces, and let $T\in\cL(X,Y)$ be a bounded
linear operator.  We write $\cK(X,Y)$ for the set of compact operators
in $\cL(X,Y)$.  We write $x_n\weakto x$ for weak convergence in $X$.
Recall that $T$ is \emph{compact} when it maps bounded subsets of $X$
into relatively norm-compact subsets of $Y$.

There are two superficially similar properties that are worth separating:
\begin{align}
  & T\text{ is compact}; \label{eq:compact}\tag{C}\\
  & x_n\weakto x\quad\Longrightarrow\quad
       \|Tx_n-Tx\|_Y\longrightarrow 0.\label{eq:seq}\tag{S}
\end{align}
In this note, property \eqref{eq:seq} is called \emph{complete continuity}.
Terminology varies: in some older literature the same words have been
used for compactness; see \cite{eom}.

The point is not to introduce a new notion.  It is to record exactly
where the implications between \eqref{eq:compact} and \eqref{eq:seq}
use information about the domain space.

\section{Compactness upgrades weak convergence}\label{sec:forward}

\begin{proposition}\label{prop:compact-to-strong}
If $T\in\cK(X,Y)$ and $x_n\weakto x$ in $X$, then
$\|Tx_n-Tx\|_Y\to 0$.
\end{proposition}

\begin{proof}
A weakly convergent sequence is norm-bounded, so compactness implies
that every subsequence of $(Tx_n)$ has a further norm-convergent
subsequence.  On the other hand, bounded linear maps preserve weak
convergence, and hence $Tx_n\weakto Tx$ in $Y$.  If a subsequence of
$(Tx_n)$ converges in norm to $y$, it converges weakly to $y$ as well;
uniqueness of the weak limit gives $y=Tx$.

If $Tx_n$ did not converge to $Tx$ in norm, some subsequence would remain
at distance at least $\varepsilon>0$ from $Tx$.  Compactness would give
a further norm-convergent subsequence, whose limit must equal $Tx$,
a contradiction.
\end{proof}

This proof contains the mechanism we will reuse in
\cref{lem:weak-cauchy}: \emph{relative norm-compactness produces limit
points, while weak information identifies them}.  Compare the treatment
of compact operators in \cite{lecture-notes}.

\section{When is the converse true?}\label{sec:converse}

The reverse implication is not a purely formal consequence of
\cref{prop:compact-to-strong}; the missing ingredient is a supply of
weakly convergent subsequences.

\begin{proposition}\label{prop:reflexive-converse}
Suppose $X$ is reflexive.  If $T\in\cL(X,Y)$ satisfies
\eqref{eq:seq}, then $T$ is compact.
\end{proposition}

\begin{proof}
Let $(x_n)$ be bounded in $X$.  Reflexivity, together with weak
sequential compactness of bounded subsets of a reflexive Banach space,
gives a subsequence $x_{n_k}\weakto x$.  By \eqref{eq:seq},
$Tx_{n_k}\to Tx$ in norm.  Thus the image under $T$ of every bounded
sequence admits a norm-convergent subsequence, which is precisely the
sequential criterion for compactness.
\end{proof}

\begin{remark}\label{rem:hinge}
The use of reflexivity occurs at exactly one point: extracting a weakly
convergent subsequence from an arbitrary bounded sequence.  This is the
assumption to revisit when comparing the proposition with the example
below, rather than treating the entire proof as an indivisible black box.
\end{remark}

\begin{example}[The obstruction in $\ell^1$]\label{ex:schur}
The space $\ellone$ has the \emph{Schur property}: a weakly convergent
sequence in $\ellone$ converges in norm.  Consequently the identity
$\mathrm{Id}_{\ellone}$ satisfies \eqref{eq:seq}.  It is not compact:
the canonical unit vectors $(e_n)$ form a bounded sequence with
$\|e_n-e_m\|_1=2$ for $n\neq m$ and hence no norm-convergent
subsequence.  This is a counterexample to the converse without an
assumption on $X$; see also \cite{eom}.
\end{example}

The juxtaposition of \cref{prop:reflexive-converse,ex:schur} suggests
that the real issue is not the terminology ``completely continuous'',
but the geometry of bounded sequences in $X$.

\section{A link through weakly Cauchy sequences}\label{sec:cauchy}

A sequence $(x_n)$ is \emph{weakly Cauchy} when $(f(x_n))$ is Cauchy for
every $f\in X^*$.  Unlike a weakly convergent sequence, a weakly Cauchy
sequence need not have a weak limit in $X$.

\begin{lemma}\label{lem:weak-cauchy}
If $T\in\cK(X,Y)$ and $(x_n)$ is weakly Cauchy in $X$, then
$(Tx_n)$ converges in norm in $Y$.
\end{lemma}

\begin{proof}
By the uniform boundedness principle, $(x_n)$ is bounded.  Thus
$(Tx_n)$ is relatively norm-compact.  For every $g\in Y^*$, the scalar
sequence $g(Tx_n)=(T^*g)(x_n)$ is Cauchy, hence has a unique limit.
Any two norm-convergent subsequences of $(Tx_n)$ must therefore have
the same limit, because functionals in $Y^*$ separate points.  Relative
norm-compactness now forces the entire sequence to converge in norm.
\end{proof}

The relation to \cref{prop:compact-to-strong} is transparent: there we
knew the weak limit in advance; here scalar Cauchy information is enough
to show that \emph{all} possible norm-limit points agree.

\paragraph{A possible line to pursue.}\label{par:possible-line}
In \cref{prop:reflexive-converse}, reflexivity was used only to extract
weakly convergent subsequences.  One might instead look for a condition
on $X$ guaranteeing weakly Cauchy subsequences and separately require
$T$ to send such sequences to norm-Cauchy sequences.  That would also
imply compactness by the same sequential argument.  This is a way of
\emph{organizing a question}, not a claim that a new theorem has been
found.  The difference between weak convergence and weak Cauchyness
must remain explicit; \cref{ex:schur} does not erase that difference.

\section{Loose ends and reading pointers}\label{sec:threads}

The current comparison leaves a specific task for later reading: which
hypotheses on $X$ make bounded sequences sufficiently well behaved for
the argument in \cref{prop:reflexive-converse} to go through?  A related
question is whether replacing the sequential property \eqref{eq:seq}
by a condition on weakly Cauchy sequences is natural in a particular
problem.  These are prompts for subsequent work, not assertions about
what is open in the literature.

When adding a paper or a handwritten calculation to this note, the
useful edit is not to create a new catalogue entry.  It is to put the
relevant statement or obstruction \emph{next to the argument it changes},
with a citation, a precise hypothesis, and a cross-reference to any
other affected passage.  If an argument is unfinished, say exactly where
it stops rather than converting it into a theorem.

\begin{thebibliography}{9}

\bibitem{eom}
\emph{Encyclopedia of Mathematics},
``Completely-continuous operator.''
\url{https://encyclopediaofmath.org/wiki/Completely-continuous_operator}

\bibitem{lecture-notes}
\emph{Lecture notes on operator ideals}, \S2.1, ``Compact operators in
Banach spaces.''
\url{https://users.fmf.uni-lj.si/kostenko/teach/IdealsNotes.pdf}

\end{thebibliography}

\end{document}
